\documentclass[11pt]{article}
\usepackage[margin=2cm]{geometry}
\usepackage{hyperref}
\hypersetup{
    colorlinks=true,
    linkcolor=blue,
    citecolor=red,
    filecolor=magenta,      
    urlcolor=blue
    }
\usepackage{multirow}
\usepackage{amssymb,amsmath}
\usepackage{mathrsfs}
\usepackage[nameinlink,capitalize]{cleveref}

\usepackage{algorithm,algpseudocodex}
\renewcommand{\vec}[1]{\boldsymbol{#1}}
\newcommand{\intomega}[1]{\langle{#1}\rangle_\Omega}
\newcommand{\adjoint}[1]{{#1}^{\text{ad}}}
\newcommand{\xir}{\mathring{\widetilde{\xi}}}
\title{Lippmann Schwinger solver for linear elasticity}
\author{Eike Mueller, University of Bath}
\begin{document}
\maketitle
%%%%%%%%%%%%%%%%%%%%%%%%%%%%%%%%%%%%%%%%%%%%%%%%%%%%%%%%%%%%%%%%%%%%%%%%
\section{Problem description}
%%%%%%%%%%%%%%%%%%%%%%%%%%%%%%%%%%%%%%%%%%%%%%%%%%%%%%%%%%%%%%%%%%%%%%%%
We aim to solve the equations of linear elasticity
\begin{equation}
    \begin{aligned}
        \partial_j \sigma_{ij} &= 0 \\
        \sigma_{ij} &= C_{ijk\ell} \varepsilon_{k\ell}\qquad\text{with $\varepsilon_{k\ell} = \varepsilon^*_{k\ell} + \overline{\varepsilon}_{k\ell}$}\\
        \varepsilon^*_{k\ell} &= \frac{1}{2}\left(\partial_k u_\ell + \partial_\ell u_k\right)
    \end{aligned}\label{eqn:continuum_equations}
\end{equation}
in the domain $\Omega=[0,L_0]\times[0,L_0]\times[0,L_1]$
Here $C=C(x)$ is the given spatially varying elasticity tensor which has the following form for an isotropic material:
\begin{equation}
    C(x) = \lambda(x) \delta_{ij}\delta_{k\ell} + \mu(x) (\delta_{ik}\delta_{j\ell} + \delta_{i\ell}\delta_{jk})
\end{equation}
$\overline{\varepsilon}$ is a constant symmetric tensor representing the average strain and $u=u(x)$ is periodic such that $\int_\Omega u_i(x)\;dx = 0$ and $\int_\Omega \partial_i u_j(x)\;dx = 0$, i.e. $\int_\Omega \varepsilon^*_{ij}(x)\;dx=0$. This implies that $\frac{1}{|\Omega|}\int_\Omega \varepsilon_{ij}(x)\;dx=\overline{\varepsilon}_{ij}$
%%%%%%%%%%%%%%%%%%%%%%%%%%%%%%%%%%%%%%%%%%%%%%%%%%%%%%%%%%%%%%%%%%%%%%%%
\section{Discretisation}
%%%%%%%%%%%%%%%%%%%%%%%%%%%%%%%%%%%%%%%%%%%%%%%%%%%%%%%%%%%%%%%%%%%%%%%%
Consider the staggered grid in Fig. 1 of \cite{gelebart2020modified}, with the placement of variables shown at the voxel centres $\mathcal{C}$ and voxel corners (vertices) $\mathcal{V}$ in \cref{tab:placement}.
\begin{table}
    \caption{Placement of variables on grid}\label{tab:placement}
    \begin{center}
\begin{tabular}{ccc}
    \hline
location & variable & \\
\hline\hline
\multirow{2}{*}{voxel centres $\mathcal{C}$} & strain & $\varepsilon^*_{ij}=\frac{1}{2}(\partial_i u_j + \partial_j u_i)$ \\
& stress & $\sigma_{ij}=C_{ijk\ell} \varepsilon_{k\ell}$ \\\hline
\multirow{2}{*}{voxel corners $\mathcal{V}$} & displacement & $u_i$ \\
& stress divergence & $\partial_i \sigma_{ij}$ \\
\hline
\end{tabular}
\end{center}
\end{table}
Define the forward gradient $\vec{D}^+:\mathcal{V}\rightarrow \mathcal{C}$ of a variable $f\in \mathcal{V}$ placed at the voxel corners as
\begin{equation}
    \begin{aligned}
(D_0^+ f)_{abc} &= \frac{1}{4h_0}\Big[\left(f_{a+1,b,c}+f_{a+1,b+1,c}+f_{a+1,b,c+1}+f_{a+1,b+1,c+1}\right)\\
& \qquad -\;\;\left(f_{a,b,c}+f_{a,b+1,c}+f_{a,b,c+1}+f_{a,b+1,c+1} \right)\Big]
    \end{aligned}
\end{equation}
with corresponding expressions for $D_1^+ f$ and $D_2^+ f$.
The backward gradient $\vec{D}^-:\mathcal{C}\rightarrow \mathcal{V}$ of a variable $g\in \mathcal{C}$ placed at the voxel centres is given by
\begin{equation}
    \begin{aligned}
(D_0^- g)_{abc} &= \frac{1}{4h_0}\Big[\left(g_{a,b,c}+g_{a,b-1,c}+g_{a,b,c-1}+g_{a,b-1,c-1}\right)\\
& \qquad -\;\;\left(g_{a-1,b,c}+g_{a-1,b-1,c}+g_{a-1,b,c-1}+g_{a-1,b-1,c-1} \right)\Big]
    \end{aligned}
\end{equation}
Observe that $-\vec{D}^+\cdot \vec{D}^- = -\vec{D}^-\cdot \vec{D}^+$ is an approximation of the Laplace operator $\Delta=\partial_j\partial_j$.
%%%%%%%%%%%%%%%%%%%%%%%%%%%%%%%%%%%%%%%%%%%%%%%%%%%%%%%%%%%%%%%%%%%%%%%%
\section{Fourier expansion}
%%%%%%%%%%%%%%%%%%%%%%%%%%%%%%%%%%%%%%%%%%%%%%%%%%%%%%%%%%%%%%%%%%%%%%%%
The Fourier expansion of a general function $f$ is given by
\begin{equation}
    \begin{aligned}
\widehat{f}_{k_0,k_1,k_2} &= \sum_{n_0=0}^{N_0-1}\sum_{n_1=0}^{N_1-1}\sum_{n_2=0}^{N_2-1} f_{n_0,n_1,n_2}\exp\left[-\frac{2\pi i k_0 n_0}{N_0}-\frac{2\pi i k_1 n_1}{N_1}-\frac{2\pi i k_2 n_2}{N_2}\right] \\
f_{n_0,n_1,n_2} &= \frac{1}{N_0N_1N_2}\sum_{k_0=0}^{N_0-1}\sum_{k_1=0}^{N_1-1}\sum_{k_2=0}^{N_2-1} \widehat{f}_{k_0,k_1,k_2}\exp\left[\frac{2\pi i k_0 n_0}{N_0}+\frac{2\pi i k_1 n_1}{N_1}+\frac{2\pi  i k_2 n_2}{N_2}\right] 
    \end{aligned}
\end{equation}
For given $\vec{k}\in[0,N_0-1]\times[0,N_1-1]\times[0,N_2-1]$ define $\vec{\xi}\in[0,2\pi]^3$ with $\xi_i = \frac{2\pi k_i}{N_i}$, then the Fourier expansion can also be written as 
    \begin{xalignat}{2}
\widehat{f}_{\vec{\xi}} &= \sum_{n_0=0}^{N_0-1}\sum_{n_1=0}^{N_1-1}\sum_{n_2=0}^{N_2-1} f_{\vec{n}}\exp\left[-i\vec{\xi}\cdot \vec{n}\right], &
f_{\vec{n}} = \frac{1}{N_0N_1N_2}\sum_{k_0=0}^{N_0-1}\sum_{k_1=0}^{N_1-1}\sum_{k_2=0}^{N_2-1} \widehat{f}_{\vec{\xi}}\exp\left[i\vec{\xi}\cdot\vec{n}\right] 
    \end{xalignat}
It is then easy to see that in Fourier-space we have
\begin{equation}
\widehat{\vec{D}}^\pm = i \widetilde{\vec{\xi}} e^{\pm i \eta}
\end{equation}
with
\begin{equation}
\begin{aligned}
\widetilde{\xi}_0 &= \frac{2}{h_0}\sin\left(\frac{\xi_0}{2}\right)\cos\left(\frac{\xi_1}{2}\right)\cos\left(\frac{\xi_2}{2}\right)\\[1ex]
\widetilde{\xi}_1 &= \frac{2}{h_1}\cos\left(\frac{\xi_0}{2}\right)\sin\left(\frac{\xi_1}{2}\right)\cos\left(\frac{\xi_2}{2}\right)\\[1ex]
\widetilde{\xi}_2 &= \frac{2}{h_2}\cos\left(\frac{\xi_0}{2}\right)\cos\left(\frac{\xi_1}{2}\right)\sin\left(\frac{\xi_2}{2}\right)
\end{aligned}
\end{equation}
and
\begin{equation}
\eta = \frac{1}{2}\left(\xi_0+\xi_1+\xi_2\right).
\end{equation}
Observe that $\widehat{D}_i^+ \widehat{D}_j^- = \widehat{D}_i^- \widehat{D}_j^-= -\widetilde{\xi}_i \widetilde{\xi}_j$.
Consider the discretised homogeneous equations
\begin{xalignat}{2}
\sigma_{ij} &= C^0_{ijk\ell} D^+_k u_\ell + \tau_{ij}, & 
D^-_j \sigma_{ij} &= 0\label{eqn:discretised_equations}
\end{xalignat}
for a reference material with constant $C^0$. In Fourier-space this becomes
\begin{xalignat}{2}
\widehat{\sigma}_{ij} &= ie^{i\eta} C^0_{ijk\ell} \widetilde{\xi}_k \widehat{u}_\ell + \widehat{\tau}_{ij}, & 
ie^{-i\eta}\widetilde{\xi}_j \widehat{\sigma}_{ij} &= 0
\end{xalignat}
and therefore
\begin{equation}
K^0_{ik} \widehat{u}_k = ie^{-i\eta} \widehat{\tau}_{ij}\widetilde{\xi}_j
\end{equation}
with the acoustic tensor
\begin{equation}
    K^0_{ij} = C^0_{ikj\ell} \widetilde{\xi}_k\widetilde{\xi}_{\ell}.
\end{equation}
We also have
\begin{equation}
\widehat{\varepsilon}^*_{k\ell} = \frac{i}{2} e^{i\eta} \left(\widetilde{\xi}_k \widehat{u}_\ell+\widetilde{\xi}_\ell \widehat{u}_k\right)
\end{equation}
Putting everything together and following the derivation in the appendix of \cite{moulinec1998numerical} we find that for a homogeneous and isotropic reference material with
\begin{equation}
C^0_{ijk\ell} = \lambda^0 \delta_{ij}\delta_{k\ell} + \mu^0 \left(\delta_{ik}\delta_{j\ell}+\delta_{i\ell}\delta_{jk}\right)
\end{equation}
we have
\begin{equation}
    \begin{aligned}
\widehat{\varepsilon}^*_{k\ell ij} &= -\widehat{\Gamma}^0_{k\ell ij} \widehat{\tau}_{ij}\\
K^0_{ij} &= (\lambda^0+\mu^0)\widetilde{\xi}_i\widetilde{\xi}_j + \mu^0 |\widetilde{\vec{\xi}}|^2 \delta_{ij}\\
N^0_{ij} =\left((K^0)^{-1}\right)_{ij} &= \frac{1}{\mu^0 |\widetilde{\vec{\xi}}|^2}\left(\delta_{ij}-\frac{\widetilde{\xi}_i\widetilde{\xi}_j}{|\widetilde{\vec{\xi}}|^2}\frac{\lambda^0+\mu^0}{\lambda^0+2\mu^0}\right)\\
\widehat{\Gamma}^0_{k\ell ij} &= \frac{1}{4}\left(N^0_{\ell i}\widetilde{\xi}_j\widetilde{\xi}_k+N^0_{k i}\widetilde{\xi}_j\widetilde{\xi}_\ell + N^0_{\ell j}\widetilde{\xi}_i\widetilde{\xi}_k+N^0_{kj}\widetilde{\xi}_i\widetilde{\xi}_\ell\right)\\
&= \frac{1}{\mu^0}\left(\frac{1}{4}\left(\delta_{\ell i} \mathring{\widetilde{\xi}}_j\mathring{\widetilde{\xi}}_k+\delta_{ki} \mathring{\widetilde{\xi}}_j\mathring{\widetilde{\xi}}_\ell+\delta_{\ell j} \mathring{\widetilde{\xi}}_i\mathring{\widetilde{\xi}}_k+\delta_{k j} \mathring{\widetilde{\xi}}_i\mathring{\widetilde{\xi}}_\ell\right)-\frac{\lambda^0+\mu^0}{\lambda^0+2\mu^0}\mathring{\widetilde{\xi}}_i\mathring{\widetilde{\xi}}_j\mathring{\widetilde{\xi}}_k\mathring{\widetilde{\xi}}_\ell\right)\\
&= \frac{1}{\mu^0}\left( \widehat{A}_{k\ell ij} - \frac{\lambda^0+\mu^0}{\lambda^0+2\mu^0} \widehat{B}_{k\ell ij}\right)
    \end{aligned}\label{eqn:fourier_expression}
\end{equation}
with
\begin{equation}
\mathring{\widetilde{\xi}}_i = \begin{cases}
0 & \text{if $\|\widetilde{\vec{\xi}}\|=0$}\\
\widetilde{\xi}_i/\|\widetilde{\vec{\xi}}\| & \text{otherwise}
\end{cases}
\end{equation}
where $\|\cdot\|$ denotes the Euclidean two-norm and we defined
\begin{equation}
\begin{aligned}
\widehat{A}_{k\ell ij} &= \frac{1}{4}\left(\delta_{\ell i} \mathring{\widetilde{\xi}}_j\mathring{\widetilde{\xi}}_k+\delta_{ki} \mathring{\widetilde{\xi}}_j\mathring{\widetilde{\xi}}_\ell+\delta_{\ell j} \mathring{\widetilde{\xi}}_i\mathring{\widetilde{\xi}}_k+\delta_{k j} \mathring{\widetilde{\xi}}_i\mathring{\widetilde{\xi}}_\ell\right)\\
\widehat{B}_{k\ell ij} &= \mathring{\widetilde{\xi}}_i\mathring{\widetilde{\xi}}_j\mathring{\widetilde{\xi}}_k\mathring{\widetilde{\xi}}_\ell
\end{aligned}
\end{equation}
If we use Voigt's notation \cite{voigt1928lehrbuch} and the convention used in the AMITEX code to arrange the six independent components of the symmetric $3\times 3$ tensors $\widehat{\varepsilon}$ and $\widehat{\tau}$ into vectors
\begin{equation}
    \begin{aligned}
    \widehat{\vec{\varepsilon}}^* &=
    \begin{pmatrix}
        \widehat{\varepsilon}^*_{0} &
        \widehat{\varepsilon}^*_{1} &
        \widehat{\varepsilon}^*_{2} &
        \widehat{\varepsilon}^*_{3} &
        \widehat{\varepsilon}^*_{4} &
        \widehat{\varepsilon}^*_{5}
    \end{pmatrix}:= 
    \begin{pmatrix}
        \widehat{\varepsilon}^*_{00} &
        \widehat{\varepsilon}^*_{11} &
        \widehat{\varepsilon}^*_{22} &
        \widehat{\varepsilon}^*_{01} &
        \widehat{\varepsilon}^*_{02} &
        \widehat{\varepsilon}^*_{12}
    \end{pmatrix}\\
    \widehat{\vec{\tau}} &= 
    \begin{pmatrix}
        \widehat{\tau}_{0} &
        \widehat{\tau}_{1} &
        \widehat{\tau}_{2} &
        \widehat{\tau}_{3} &
        \widehat{\tau}_{4} &
        \widehat{\tau}_{5}
    \end{pmatrix}= 
    \begin{pmatrix}
        \widehat{\tau}_{00} &
        \widehat{\tau}_{11} &
        \widehat{\tau}_{22} &
        \widehat{\tau}_{01} &
        \widehat{\tau}_{02} &
        \widehat{\tau}_{12}
    \end{pmatrix}
\end{aligned}
\end{equation}
we can write
\begin{equation}
\widehat{\vec{\varepsilon}}^* = -\frac{1}{\mu^0}\left(\widehat{\vec{\varepsilon}}^{*(A)} - \frac{\lambda^0+\mu^0}{\lambda^0+2\mu^0} \widehat{\vec{\varepsilon}}^{*(B)}\right)
\end{equation}
with 
\begin{equation}
\begin{aligned}    
    \widehat{\vec{\varepsilon}}^{*(A)} &= \begin{pmatrix}
    \xir_0^2 \widehat{\tau}_0 + \xir_0(\xir_2\widehat{\tau}_4+\xir_1\widehat{\tau}_3)\\[2ex]
    \xir_1^2 \widehat{\tau}_1 + \xir_1(\xir_2\widehat{\tau}_5+\xir_0\widehat{\tau}_3)\\[2ex]
    \xir_2^2 \widehat{\tau}_2 + \xir_2(\xir_1\widehat{\tau}_5+\xir_0\widehat{\tau}_4)\\[2ex]
    \frac{1}{2}\left(\xir_0\xir_1(\widehat{\tau}_0+\widehat{\tau}_1)+(\xir_0^2+\xir_1^2)\widehat{\tau}_3+\xir_2(\xir_0\widehat{\tau}_5+\xir_1\widehat{\tau}_4)\right)\\[2ex]
    \frac{1}{2}\left(\xir_0\xir_2(\widehat{\tau}_0+\widehat{\tau}_2)+(\xir_0^2+\xir_2^2)\widehat{\tau}_4+\xir_1(\xir_0\widehat{\tau}_5+\xir_2\widehat{\tau}_3)\right)\\[2ex]
    \frac{1}{2}\left(\xir_1\xir_2(\widehat{\tau}_1+\widehat{\tau}_2)+(\xir_1^2+\xir_2^2)\widehat{\tau}_5+\xir_0(\xir_1\widehat{\tau}_4+\xir_2\widehat{\tau}_3)\right)
    \end{pmatrix}
\end{aligned}
\end{equation}
and
\begin{equation}
\begin{aligned}    
    \widehat{\vec{\varepsilon}}^{*(B)} &= 
    \left(
    \xir_0^2\widehat{\tau}_0+ \xir_1^2\widehat{\tau}_1+ \xir_2^2\widehat{\tau}_2+ 2(\xir_0\xir_1\widehat{\tau}_3+ \xir_0 \xir_2\widehat{\tau}_4+\xir_1\xir_2\widehat{\tau}_5)
    \right)\begin{pmatrix}
    \xir_0^2 \\[2ex] \xir_1^2 \\[2ex] \xir_2^2 \\[2ex] \xir_0\xir_1 \\[2ex] \xir_0 \xir_2 \\[2ex] \xir_1\xir_2
    \end{pmatrix}^\top 
\end{aligned}
\end{equation}
where the expression for $\widehat{\vec{\varepsilon}}^{*(A)}=\widehat{A}\widehat{\vec{\tau}}$ follows from the $6\times 6$ matrix representation of $\widehat{A}_{k\ell ij}$ in Voigt's notation
\begin{equation}
\widehat{A} = \begin{pmatrix}
    \xir_0^2 & 0 & 0 & \xir_0\xir_1 & \xir_0\xir_2 & 0 \\[2ex]
    0 & \xir_1^2 & 0 & \xir_0\xir_1 & 0 & \xir_1\xir_2\\[2ex]
    0 & 0 & \xir_2^2 & 0 & \xir_0\xir_2 & \xir_1\xir_2\\[2ex]
    \frac{1}{2}\xir_0\xir_1 & \frac{1}{2}\xir_0\xir_1 & 0 & \frac{1}{2}(\xir_0^2+\xir_1^2) & \frac{1}{2}\xir_1\xir_2 & \frac{1}{2}\xir_0\xir_2\\[2ex]
    \frac{1}{2}\xir_0\xir_2 & 0 & \frac{1}{2}\xir_0\xir_2 & \frac{1}{2}\xir_1\xir_2 & \frac{1}{2}(\xir_0^2+\xir_2^2) & \frac{1}{2}\xir_0\xir_1\\[2ex]
    0 & \frac{1}{2}\xir_1\xir_2 & \frac{1}{2}\xir_1\xir_2 & \frac{1}{2}\xir_0\xir_2 & \frac{1}{2}\xir_0\xir_1 & \frac{1}{2}(\xir_1^2+\xir_2^2)
    \end{pmatrix}
\end{equation}
%%%%%%%%%%%%%%%%%%%%%%%%%%%%%%%%%%%%%%%%%%%%%%%%%%%%%%%%%%%%%%%%%%%%%%%%
\section{Lippmann Schwinger iteration}
%%%%%%%%%%%%%%%%%%%%%%%%%%%%%%%%%%%%%%%%%%%%%%%%%%%%%%%%%%%%%%%%%%%%%%%%
To solve \eqref{eqn:continuum_equations}, observe that the discretisation of these equations is given by \eqref{eqn:discretised_equations} with
\begin{equation}
    \begin{aligned}
    \tau_{ij} &= (C_{ijk\ell} - C^0_{ijk\ell}) \varepsilon_{k\ell} \\
    &= (\lambda-\lambda^0) \varepsilon_{kk} \delta_{ij} + 2(\mu-\mu^0)\varepsilon_{ij}.
    \end{aligned}
\end{equation}
We can write $\tau_{ij} = C_{ijk\ell}  \varepsilon_{k\ell}$ explicitly as
\begin{equation}
\begin{aligned}
\tau_{00} &= 2\mu\varepsilon_{00} + \lambda\left(\varepsilon_{00}+\varepsilon_{11}+\varepsilon_{22}\right) \\
\tau_{11} &= 2\mu\varepsilon_{11} + \lambda\left(\varepsilon_{00}+\varepsilon_{11}+\varepsilon_{22}\right)\\
\tau_{22} &= 2\mu\varepsilon_{22} + \lambda\left(\varepsilon_{00}+\varepsilon_{11}+\varepsilon_{22}\right)\\
\tau_{01} &= 2\mu\varepsilon_{01}\\
\tau_{02} &= 2\mu\varepsilon_{02}\\
\tau_{12} &= 2\mu\varepsilon_{12}.
\end{aligned}\label{eqn:tau_computation}
\end{equation}
This leads to the Lippmann Schwinger iteration shown in \cref{alg:lippmann_schwinger}. Alternatively the method can be written in incremental form as in \cref{alg:lippmann_schwinger_incremental}; this uses the fact that $\Gamma^0*C^0 \varepsilon=\varepsilon$.
\begin{algorithm}
    \caption{Lippmann Schwinger iteration: simplest formulation}
    \label{alg:lippmann_schwinger}
\begin{algorithmic}[1]
    \State{Set $\varepsilon=\overline{\varepsilon}$}
    \For{$n=0,1,\dots$}
    \State{Compute stress $\sigma_{ij}=C_{ijk\ell}\varepsilon_{k\ell}$ according to \eqref{eqn:tau_computation}}
    \State{Check convergence based on $D^-_j \sigma_{ij}$}
    \State{Compute stress correction $\tau_{ij}=\sigma_{ij}-C^0_{ijk\ell} \varepsilon_{k\ell}$}
    \State{Compute $\widehat{\tau}_{ij} = \mathscr{F}[\tau_{ij}]$}\Comment{Fourier-transform}
    \State{Set $\widehat{\varepsilon}^*_{k\ell} = -\widehat{\Gamma}^0_{k\ell ij} \widehat{\tau}_{ij}$}\Comment{Solve in Fourier space}
    \State{Update $\varepsilon_{k\ell} = \mathscr{F}^{-1}[\widehat{\varepsilon}^*_{k\ell}]+\overline{\varepsilon}_{k\ell}$}\Comment{Inverse Fourier-transform}        
    \EndFor
\State{\Return $\varepsilon$}
\end{algorithmic}
\end{algorithm}
\begin{algorithm}
    \caption{Lippmann Schwinger iteration: incremental formulation}
    \label{alg:lippmann_schwinger_incremental}
\begin{algorithmic}[1]
    \State{Set $\varepsilon=\overline{\varepsilon}$}
    \For{$n=0,1,\dots$}
    \State{Compute stress $\sigma_{ij}=C_{ijk\ell}\varepsilon_{k\ell}$ according to \eqref{eqn:tau_computation}}
    \State{Check convergence based on $D^-_j \sigma_{ij}$}    
    \State{Compute $\widehat{\sigma}_{ij} = \mathscr{F}[\sigma_{ij}]$}\Comment{Fourier-transform}
    \State{Compute $\widehat{r}_{k\ell} = -\widehat{\Gamma}^0_{k\ell ij} \widehat{\sigma}_{ij}$}\Comment{Solve in Fourier space}
    \State{Compute $r_{ij} = \mathscr{F}^{-1}[\widehat{r}_{ij}]$}\Comment{Inverse Fourier-transform}
    \State{Update $\varepsilon_{k\ell} \mapsto \varepsilon_{k\ell} + r_{k\ell}$}
    \EndFor
\State{\Return $\varepsilon$}
\end{algorithmic}
\end{algorithm}
\subsection{Stopping criterion}
The stopping criterion in \cite{moulinec1998numerical} is given by
\begin{equation}
    R(\sigma) = \frac{\sqrt{\langle \|\nabla \cdot \sigma\|^2\rangle_\Omega}}{\|\langle\sigma\rangle_\Omega\|} \le \epsilon\label{eqn:stopping_criterion}
\end{equation}
where $\|\cdot\|$ denotes the two-norm and $\langle f \rangle_\Omega = \frac{1}{|\Omega|}\int_\Omega f\;dx$ is the spatial average of $f$. In a discrete setting the left hand side of \eqref{eqn:stopping_criterion} can be written as
\begin{equation}
    R(\sigma) = \frac{\sqrt{\langle\|D^-\cdot \sigma\|^2\rangle}}{\|\langle \sigma \rangle\|} = \frac{\sqrt{N\langle\|\xi\cdot \widehat{\sigma}\|^2\rangle}}{\|\widehat{\sigma}_{\vec{\xi}=0}\|} =: \widehat{R}(\widehat{\sigma})
\end{equation}
where $\widehat{\sigma}$ is the Fourier transform of $\sigma$ and now $\langle f \rangle = \frac{1}{N}\sum_{\vec{n}} f_{\vec{n}}$ with $N:=N_0N_1N_2$ the total number of voxels.
Define
\begin{equation}
    \begin{aligned}
\delta &:= D^-\cdot \sigma = 
\begin{pmatrix}
    D^-_0 \sigma_0 + D^-_1 \sigma_3 + D^-_2 \sigma_4\\
    D^-_0 \sigma_3 + D^-_1 \sigma_1 + D^-_2 \sigma_5\\
    D^-_0 \sigma_4 + D^-_1 \sigma_5 + D^-_2 \sigma_2
\end{pmatrix}
\\[2ex]
\widehat{\delta} &:= \xi \cdot \widehat{\sigma} =
\begin{pmatrix}
    \xi_0 \widehat{\sigma}_0 + \xi_1 \widehat{\sigma}_3 + \xi_2 \widehat{\sigma}_4\\
    \xi_0 \widehat{\sigma}_3 + \xi_1 \widehat{\sigma}_1 + \xi_2 \widehat{\sigma}_5\\
    \xi_0 \widehat{\sigma}_4 + \xi_1 \widehat{\sigma}_5 + \xi_2 \widehat{\sigma}_2
\end{pmatrix}
\end{aligned}
\end{equation}
to obtain
\begin{equation}
R(\sigma) = \frac{\|\delta\|}{\sqrt{N}\|\overline{\sigma}\|}
\end{equation}
where
\begin{xalignat}{3}
    \|\delta\| &= \sqrt{\sum_i\sum_{\vec{n}}\delta^2_{i\vec{n}}}, &
    \overline{\sigma_i} &= \frac{1}{N}\sum_{\vec{n}}\sigma_{i\vec{n}}, &
    \|\overline{\sigma}\| &= \sqrt{\overline{\sigma}_0^2+\overline{\sigma}_1^2+\overline{\sigma}_2^2 + 2(\overline{\sigma}_3^2+\overline{\sigma}_4^2+\overline{\sigma}_5^2)}
\end{xalignat}
Similarly, we find 
\begin{equation}
\widehat{R}(\widehat{\sigma}) = \frac{\|\widehat{\delta}\|}{\|\widehat{\sigma}_{\vec{\xi}=0}\|}
\qquad\text{with\quad $\|\widehat{\delta}\| = \sqrt{\sum_i\sum_{\vec{\xi}}|\widehat{\delta}_{i\vec{\xi}}|^2}$}.
\end{equation}
\subsection{Anderson acceleration}
The Lippmann-Schwinger algorithm with Anderson acceleration \cite{wicht2021anderson} is shown in \cref{alg:lippmann_schwinger_anderson}. It requires additional storage of $2\times(d+1)$ tensors for the state $\varepsilon^{(s)}$ and residuals $r^{(s)}$, as well as a $(d+1)\times (d+1)$ matrix and vectors $u,v\in\mathbb{R}^{d+1}$.
\begin{algorithm}
    \caption{Lippmann Schwinger iteration: Anderson acceleration with depth $d$}\label{alg:lippmann_schwinger_anderson}
\begin{algorithmic}[1]
    \State{Set $\varepsilon=E$}
    \State{Initialise $\varepsilon^{(0)} = \varepsilon$, $\varepsilon^{(j)} = 0$ for $j=1,2,\dots,d$}
    \State{Initialise $r^{(0)} = 0$ for $s=0,1,2,\dots,d$}
    \State{Initialise $(d+1)\times (d+1)$ identity matrix $A=\mathbb{I}$}
    \For{$n=0,1,\dots$}
    \State{Compute stress $\sigma_{ij}=C_{ijk\ell}\varepsilon_{k\ell}$ according to \eqref{eqn:tau_computation}}
    \State{Check convergence based on $D^-_j \sigma_{ij}$}    
    \State{Compute $\widehat{\sigma}_{ij} = \mathscr{F}[\sigma_{ij}]$}\Comment{Fourier-transform}
    \State{Compute $\widehat{r}_{k\ell} = \widehat{\Gamma}^0_{k\ell ij} \widehat{\sigma}_{ij}$}\Comment{Solve in Fourier space}
    \State{Set $m_k=\min\{n,d\}$}
    \For{$s=m_k,m_k-1,\dots,2,1$}
    \State{Set $r^{(s)}\gets r^{(s-1)}$}
    \For{$t=m_k,m_k-1,\dots,2,1$}
    \State{Set $A_{st}\gets A_{s-1,t-1}$}
    \EndFor
    \EndFor
    \State{Compute $r^{(0)}_{ij} = \mathscr{F}^{-1}[\widehat{r}_{ij}]$}\Comment{Inverse Fourier-transform}
    \For{$s=0,1,2,\dots,m_k$}
    \State{Set $A_{0s} = A_{s0}\gets r^{(0)}_{ij}r^{(s)}_{ij}$}
    \EndFor
    \State{Solve $Av=u$ with $u=(1,1,\dots,1)\in\mathbb{R}^d$}
    \State{Set $\alpha = v/(u^\top v)$}
    \State{Set $\varepsilon_{ij} = \sum_{s=0}^{d} \alpha_s (\varepsilon^{(s)}_{ij}-r^{(s)}_{ij})$}
    \For{$s=m_k,m_k-1,\dots,2,1$}
    \State{Set $\varepsilon^{(s)}\gets \varepsilon^{(s-1)}$}
    \EndFor
    \State{Set $\varepsilon_{ij} \gets \varepsilon_{ij}^{(0)}$}
    \EndFor
\State{\Return $\varepsilon$}
\end{algorithmic}
\end{algorithm}
%%%%%%%%%%%%%%%%%%%%%%%%%%%%%%%%%%%%%%%%%%%%%%%%%%%%%%%%%%%%%%%%%%%%%%%
\section{Anisotropic case}
%%%%%%%%%%%%%%%%%%%%%%%%%%%%%%%%%%%%%%%%%%%%%%%%%%%%%%%%%%%%%%%%%%%%%%%
In this case, we have that
\begin{equation}
\sigma_{ij} = C_{ijk\ell}\varepsilon_{k\ell}\label{eqn:stress_strain_anisotropic}
\end{equation}
where due to the symmetries of $C$ only $21$ components of the $3\times 3\times 3 \times 3$ tensor $C$ are independent. Using Voigt-notation these can be taken to be the $21$ independent entries $C_{i}$ of a symmetric $6\times 6$ matrix $C$ defined by
\input{generated/anisotropic_elasticity.tex}
Observe that in the isotropic case the only non-zero entries $C_j$ are
\begin{equation}
\begin{aligned}
C_0 = C_1 = C_2 &= 2\mu + \lambda,\\
C_3 = C_4 = C_5 &= \mu,\\
C_6 = C_7 = C_8 &= \lambda.
\end{aligned}
\end{equation}
In Voigt-notation \eqref{eqn:stress_strain_anisotropic} becomes:
\input{generated/anisotropic_stress_strain.tex}
The tensor $\widehat{\Gamma}^{0}_{k\ell ij}$ also needs to be computed slightly differently. Let the elasticity tensor of the reference material be $C^0$; again this has $21$ independent components.
For each Fourier mode compute the symmetric $3\times 3$ matrix $K^{0}$ with
% ---- acoustic matrix K^0 ----
\begin{equation}\begin{aligned}
K^0_{00} &= {C^{0}_{0}} {{\mathring{{\widetilde{{\xi}}}}}_{0}}^{2} + {C^{0}_{3}} {{\mathring{{\widetilde{{\xi}}}}}_{1}}^{2} + {C^{0}_{4}} {{\mathring{{\widetilde{{\xi}}}}}_{2}}^{2} + 2 {C^{0}_{9}} {{\mathring{{\widetilde{{\xi}}}}}_{0}} {{\mathring{{\widetilde{{\xi}}}}}_{1}} + 2 {C^{0}_{10}} {{\mathring{{\widetilde{{\xi}}}}}_{0}} {{\mathring{{\widetilde{{\xi}}}}}_{2}} + 2 {C^{0}_{18}} {{\mathring{{\widetilde{{\xi}}}}}_{1}} {{\mathring{{\widetilde{{\xi}}}}}_{2}}\\
K^0_{10} &= {C^{0}_{3}} {{\mathring{{\widetilde{{\xi}}}}}_{0}} {{\mathring{{\widetilde{{\xi}}}}}_{1}} + {C^{0}_{6}} {{\mathring{{\widetilde{{\xi}}}}}_{0}} {{\mathring{{\widetilde{{\xi}}}}}_{1}} + {C^{0}_{9}} {{\mathring{{\widetilde{{\xi}}}}}_{0}}^{2} + {C^{0}_{11}} {{\mathring{{\widetilde{{\xi}}}}}_{0}} {{\mathring{{\widetilde{{\xi}}}}}_{2}} + {C^{0}_{12}} {{\mathring{{\widetilde{{\xi}}}}}_{1}}^{2} + {C^{0}_{13}} {{\mathring{{\widetilde{{\xi}}}}}_{1}} {{\mathring{{\widetilde{{\xi}}}}}_{2}} + {C^{0}_{18}} {{\mathring{{\widetilde{{\xi}}}}}_{0}} {{\mathring{{\widetilde{{\xi}}}}}_{2}} + {C^{0}_{19}} {{\mathring{{\widetilde{{\xi}}}}}_{1}} {{\mathring{{\widetilde{{\xi}}}}}_{2}} + {C^{0}_{20}} {{\mathring{{\widetilde{{\xi}}}}}_{2}}^{2}\\
K^0_{11} &= {C^{0}_{1}} {{\mathring{{\widetilde{{\xi}}}}}_{1}}^{2} + {C^{0}_{3}} {{\mathring{{\widetilde{{\xi}}}}}_{0}}^{2} + {C^{0}_{5}} {{\mathring{{\widetilde{{\xi}}}}}_{2}}^{2} + 2 {C^{0}_{12}} {{\mathring{{\widetilde{{\xi}}}}}_{0}} {{\mathring{{\widetilde{{\xi}}}}}_{1}} + 2 {C^{0}_{14}} {{\mathring{{\widetilde{{\xi}}}}}_{1}} {{\mathring{{\widetilde{{\xi}}}}}_{2}} + 2 {C^{0}_{19}} {{\mathring{{\widetilde{{\xi}}}}}_{0}} {{\mathring{{\widetilde{{\xi}}}}}_{2}}\\
K^0_{20} &= {C^{0}_{4}} {{\mathring{{\widetilde{{\xi}}}}}_{0}} {{\mathring{{\widetilde{{\xi}}}}}_{2}} + {C^{0}_{7}} {{\mathring{{\widetilde{{\xi}}}}}_{0}} {{\mathring{{\widetilde{{\xi}}}}}_{2}} + {C^{0}_{10}} {{\mathring{{\widetilde{{\xi}}}}}_{0}}^{2} + {C^{0}_{11}} {{\mathring{{\widetilde{{\xi}}}}}_{0}} {{\mathring{{\widetilde{{\xi}}}}}_{1}} + {C^{0}_{15}} {{\mathring{{\widetilde{{\xi}}}}}_{1}} {{\mathring{{\widetilde{{\xi}}}}}_{2}} + {C^{0}_{16}} {{\mathring{{\widetilde{{\xi}}}}}_{2}}^{2} + {C^{0}_{18}} {{\mathring{{\widetilde{{\xi}}}}}_{0}} {{\mathring{{\widetilde{{\xi}}}}}_{1}} + {C^{0}_{19}} {{\mathring{{\widetilde{{\xi}}}}}_{1}}^{2} + {C^{0}_{20}} {{\mathring{{\widetilde{{\xi}}}}}_{1}} {{\mathring{{\widetilde{{\xi}}}}}_{2}}\\
K^0_{21} &= {C^{0}_{5}} {{\mathring{{\widetilde{{\xi}}}}}_{1}} {{\mathring{{\widetilde{{\xi}}}}}_{2}} + {C^{0}_{8}} {{\mathring{{\widetilde{{\xi}}}}}_{1}} {{\mathring{{\widetilde{{\xi}}}}}_{2}} + {C^{0}_{13}} {{\mathring{{\widetilde{{\xi}}}}}_{0}} {{\mathring{{\widetilde{{\xi}}}}}_{1}} + {C^{0}_{14}} {{\mathring{{\widetilde{{\xi}}}}}_{1}}^{2} + {C^{0}_{15}} {{\mathring{{\widetilde{{\xi}}}}}_{0}} {{\mathring{{\widetilde{{\xi}}}}}_{2}} + {C^{0}_{17}} {{\mathring{{\widetilde{{\xi}}}}}_{2}}^{2} + {C^{0}_{18}} {{\mathring{{\widetilde{{\xi}}}}}_{0}}^{2} + {C^{0}_{19}} {{\mathring{{\widetilde{{\xi}}}}}_{0}} {{\mathring{{\widetilde{{\xi}}}}}_{1}} + {C^{0}_{20}} {{\mathring{{\widetilde{{\xi}}}}}_{0}} {{\mathring{{\widetilde{{\xi}}}}}_{2}}\\
K^0_{22} &= {C^{0}_{2}} {{\mathring{{\widetilde{{\xi}}}}}_{2}}^{2} + {C^{0}_{4}} {{\mathring{{\widetilde{{\xi}}}}}_{0}}^{2} + {C^{0}_{5}} {{\mathring{{\widetilde{{\xi}}}}}_{1}}^{2} + 2 {C^{0}_{16}} {{\mathring{{\widetilde{{\xi}}}}}_{0}} {{\mathring{{\widetilde{{\xi}}}}}_{2}} + 2 {C^{0}_{17}} {{\mathring{{\widetilde{{\xi}}}}}_{1}} {{\mathring{{\widetilde{{\xi}}}}}_{2}} + 2 {C^{0}_{20}} {{\mathring{{\widetilde{{\xi}}}}}_{0}} {{\mathring{{\widetilde{{\xi}}}}}_{1}}\\
\end{aligned}\end{equation}


As a sanity check, we can compute $K^0$ for the isotropic case, i.e. $C^0_0=C^0_1=C^0_2=2\mu^0+\lambda^0$, $C^0_3=C^0_4=C^5_0=\mu^0$, $C^0_6=C^0_7=C^0_8=\lambda^0$. In this case we find
% ---- acoustic matrix K^0 [isotropic material] ----
\begin{equation}\begin{aligned}
K^0_{00} &= \mu^{0} {{\mathring{{\widetilde{{\xi}}}}}_{1}}^{2} + \mu^{0} {{\mathring{{\widetilde{{\xi}}}}}_{2}}^{2} + \left(\lambda^{0} + 2 \mu^{0}\right) {{\mathring{{\widetilde{{\xi}}}}}_{0}}^{2}\\
K^0_{10} &= \left(\lambda^{0} + \mu^{0}\right) {{\mathring{{\widetilde{{\xi}}}}}_{0}} {{\mathring{{\widetilde{{\xi}}}}}_{1}}\\
K^0_{11} &= \mu^{0} {{\mathring{{\widetilde{{\xi}}}}}_{0}}^{2} + \mu^{0} {{\mathring{{\widetilde{{\xi}}}}}_{2}}^{2} + \left(\lambda^{0} + 2 \mu^{0}\right) {{\mathring{{\widetilde{{\xi}}}}}_{1}}^{2}\\
K^0_{20} &= \left(\lambda^{0} + \mu^{0}\right) {{\mathring{{\widetilde{{\xi}}}}}_{0}} {{\mathring{{\widetilde{{\xi}}}}}_{2}}\\
K^0_{21} &= \left(\lambda^{0} + \mu^{0}\right) {{\mathring{{\widetilde{{\xi}}}}}_{1}} {{\mathring{{\widetilde{{\xi}}}}}_{2}}\\
K^0_{22} &= \mu^{0} {{\mathring{{\widetilde{{\xi}}}}}_{0}}^{2} + \mu^{0} {{\mathring{{\widetilde{{\xi}}}}}_{1}}^{2} + \left(\lambda^{0} + 2 \mu^{0}\right) {{\mathring{{\widetilde{{\xi}}}}}_{2}}^{2}\\
\end{aligned}\end{equation}


which is equivalent to $K^0_{ij} = (\lambda^0+\mu^0)\mathring{\widetilde{\xi}}_i\mathring{\widetilde{\xi}}_j + \mu^0 |\mathring{\widetilde{\vec{\xi}}}|^2 \delta_{ij}$.

The inverse of $K^0$ is the $3\times 3$ matrix $N^{0}$. This requires storage of additional 6 numbers for each Fourier mode. From this, build the symmetric $6\times 6$ matrix which is defined by
\input{generated/anisotropic_fourier_matrix.tex}
Finally, compute
\begin{equation}
\widehat{\sigma}_i = -\widehat{\Gamma}^{0}_{ij} \widehat{\tau}_j.
\end{equation}
%%%%%%%%%%%%%%%%%%%%%%%%%%%%%%%%%%%%%%%%%%%%%%%%%%%%%%%%%%%%%%%%%%%%%%%
\section{Reverse mode differentiation with the adjoint method}
%%%%%%%%%%%%%%%%%%%%%%%%%%%%%%%%%%%%%%%%%%%%%%%%%%%%%%%%%%%%%%%%%%%%%%%
Assume that we have an objective function
\begin{equation}
    J = J(\varepsilon(\cdot),\sigma(\cdot)).
\end{equation}
Define the \textit{partial} functional derivatives of $J$ with respect to strain $\varepsilon(x)$ and stress $\sigma(x) = C(x)\varepsilon(x)$ as
\begin{xalignat}{2}
    E(x) &:= \frac{\delta J}{\delta \varepsilon(x)}, &
    S(x) &:= \frac{\delta J}{\delta \sigma(x)},\label{eqn:J_derivatives}
\end{xalignat}
where
\begin{equation}
A:B=\sum_{ij}A_{ij}B_{ij}.\label{eqn:scalar_product}
\end{equation}
In a slight abuse of notation, here (and in the following) we implicitly assume that the quantities defined this way are in fact the Riesz-representers with respect to the scalar product in \eqref{eqn:scalar_product}:
\begin{equation}
\frac{\delta J}{\delta \varepsilon(x)}(H) = E(x):H(x) \qquad\text{for all $H(x)$}.
\label{eqn:riesz_representer}
\end{equation}
This has implications for the implementation which are discussed in \cref{sec:voigt_scalar_product}. Given $E(x)$, $S(x)$ we want to compute
\begin{equation}
\frac{\delta J}{\delta \overline{\varepsilon}},\qquad\frac{\delta J}{\delta C(x)}\label{eqn:func_deriv_anisotropic}
\end{equation}
in the general, anisotropic case and 
\begin{equation}
\frac{\delta J}{\delta \overline{\varepsilon}(x)},\qquad\frac{\delta J}{\delta \mu(x)},\qquad\frac{\delta J}{\delta \lambda(x)} \label{eqn:func_deriv_isotropic}
\end{equation}
in the isotropic case.
%%%%%%%%%%%%%%%%%%%%%%%%%%%%%%%%%%%%%%%%%%%%%%%%%%%%%%%%%%%%%%%%%%%%%%%
\subsection{Adjoint equation}
%%%%%%%%%%%%%%%%%%%%%%%%%%%%%%%%%%%%%%%%%%%%%%%%%%%%%%%%%%%%%%%%%%%%%%%
To express the derivatives in \eqref{eqn:func_deriv_anisotropic}, \eqref{eqn:func_deriv_isotropic} in terms of the given $E(x)$, $S(x)$ with the adjoint method, first note that the strain field satisfies the Lippmann-Schwinger equation
\begin{equation}
\varepsilon(x) = \overline{\varepsilon} - \int_\Omega \Gamma^0(x-y) \tau(y)\;dy\label{eqn:lippmann_schwinger}
\end{equation}
with $\tau(x) = \left(C(x)-C^0\right)\varepsilon(x)$. This can be written in condensed form as $\varepsilon = \overline{\varepsilon}-\Gamma^0 * \left((C-C^0)\varepsilon\right)$. Define the residual operator
\begin{equation}
    \mathcal{A}(\varepsilon;C,\overline{\varepsilon}) := \varepsilon + \Gamma^0 * \left((C-C^0) \varepsilon\right) - \overline{\varepsilon}\label{eqn:residual_operator}
\end{equation}
and introduce the Lagrangian
\begin{equation}
\mathcal{L}(\varepsilon,\Lambda;C,\overline{\varepsilon}) = \mathcal{J}(\varepsilon) + \int_\Omega \Lambda(z) : \mathcal{A}(\varepsilon;C,\overline{\varepsilon})\;dz\qquad\text{with $\mathcal{J}(\varepsilon) = J(\varepsilon,\sigma(\varepsilon))$}\label{eqn:lagrangian}
\end{equation}
The Lagrange multiplier $\Lambda(x)$ has been introduced to make $\mathcal{L}$ independent of $\varepsilon(x)$, i.e. we require that the total derivative with respect to $\varepsilon(x)$ vanishes:
\begin{equation}
\frac{\delta\mathcal{L}}{\delta\varepsilon(x)} = 0.
\end{equation}
To compute the derivative of the integral on the right hand side of \eqref{eqn:lagrangian}, observe that
\begin{equation}
    \begin{aligned}
        \frac{\delta}{\delta\varepsilon(x)}\int_\Omega \Lambda(z) : \mathcal{A}(\varepsilon;C,\overline{\varepsilon})(z)\;dz&=\frac{\delta}{\delta\varepsilon(x)}\left( \int_\Omega \Lambda(z):\varepsilon(z)\;dz+\int_{\Omega}\int_{\Omega}\Lambda(z) : \Gamma^0(z-y)\left(C(y)-C^0\right)\varepsilon(y)\;dy\;dz\right)\\
&=\Lambda(x) + \int_\Omega \Lambda(z):\Gamma^0(z-x)\left(C(x)-C^0\right)\;dz \\
&=\Lambda(x) +  \left(C(x)-C^0\right) \int_\Omega  \Gamma^0(x-z) \Lambda(z)\;dz
    \end{aligned}
\end{equation}
where the last identity follows from $\Gamma^0(-x) = \Gamma^0(x)$, $\Gamma^0_{ijk\ell} = \Gamma^0_{k\ell ij}$ and $C_{ijk\ell}=C_{k\ell ij}$. This leads to
\begin{equation}
0 = \frac{\delta\mathcal{L}}{\delta\varepsilon} = \frac{\delta \mathcal{J}}{\delta \varepsilon} + \Lambda + \left(C-C^0\right) \left(\Gamma^0 * \Lambda\right)
\end{equation}
Note that the derivatives in \eqref{eqn:J_derivatives} are \textit{partial} derivatives. Using the definition of the stress $\sigma(x) = C(x)\varepsilon(x)$, the \textit{total} derivative which appears on the right hand side of the adjoint equation is 
\begin{equation}
    \begin{aligned}
\frac{\delta \mathcal{J}}{\delta \varepsilon(x)} &= \frac{\delta J}{\delta \varepsilon(x)}+ C(x)\frac{\delta J}{\delta \sigma(x)}\\
&=E(x)+ C(x)S(x).
    \end{aligned}
\end{equation}
Hence, the adjoint equation is
\begin{equation}
    \Lambda + \left(C-C^0\right) : \left(\Gamma^0 * \Lambda\right) = \mathcal{E}\qquad\text{with $\mathcal{E}(x):= -\left(E(x)+ C(x)S(x)\right)$}\label{eqn:adjoint_equation}
\end{equation}
The Lippmann-Schwinger iteration for solving the adjoint equation in \eqref{eqn:adjoint_equation} is shown in \cref{alg:lippmann_schwinger_adjoint}.
\begin{algorithm}
    \caption{Adjoint Lippmann Schwinger iteration (simplest form)}
    \label{alg:lippmann_schwinger_adjoint}
\begin{algorithmic}[1]
    \State{Set $\Lambda=0$}
    \For{$n=0,1,\dots$}
    \State{Compute $\widehat{\Lambda}_{ij} = \mathscr{F}[\Lambda_{ij}]$}\Comment{Fourier-transform}
    \State{Set $\widehat{\Theta}_{k\ell} = -\widehat{\Gamma}^0_{k\ell ij} \widehat{\Lambda}_{ij}$}\Comment{Solve in Fourier space}
    \State{Compute $\Theta_{k\ell} = \mathscr{F}^{-1}[\widehat{\Theta}_{k\ell}]$}\Comment{Inverse Fourier-transform}        
    \State{Update $\Lambda_{ij}=\mathcal{E}_{ij}+(C-C^0)_{ijk\ell}\Theta_{k\ell}$}
    \State{Check convergence}
    \EndFor
\State{\Return $\Lambda$}
\end{algorithmic}
\end{algorithm}
Comparing \cref{alg:lippmann_schwinger} and \cref{alg:lippmann_schwinger_adjoint}, we see that the primal- and adjoint Lippmann Schwinger iteration only differ in the order in which $\Gamma^0$ and $C$ are applied and in the update step. Also, the algorithm is readily applied to both the isotropic and anisotropic case by use the appropriate way of computing $\Delta = C\Theta$.
%%%%%%%%%%%%%%%%%%%%%%%%%%%%%%%%%%%%%%%%%%%%%%%%%%%%%%%%%%%%%%%%%%%%%%%
\subsection{Derivatives with respect to $C(x)$ and $\overline{\varepsilon}$}
%%%%%%%%%%%%%%%%%%%%%%%%%%%%%%%%%%%%%%%%%%%%%%%%%%%%%%%%%%%%%%%%%%%%%%%
The derivative with respect to $C$ is made up of two contributions. First, the direct contribution is
\begin{equation}
\frac{\delta \mathcal{J}}{\delta C_{ijk\ell}(x)} = S_{ij}(x)\varepsilon_{k\ell}(x).
\end{equation}
Second, the adjoint contribution is obtained by observing that
\begin{equation}
    \begin{aligned}
    \frac{\delta }{\delta C_{ijk\ell}(x)} \int_\Omega \Lambda(z) : \mathcal{A}(\varepsilon;C,\overline{\varepsilon})\;dz &=\frac{\delta }{\delta C_{ijk\ell}(x)} \int_\Omega\int_\Omega \Lambda_{ab}(z)\Gamma^0_{abcd}(z-y)\left(C(y)-C^0\right)_{cdrs}\varepsilon_{rs}(y)\;dy\;dz\\
    &= \int_\Omega \Lambda_{ab}(z)\Gamma^0_{abij}(z-x)\varepsilon_{k\ell}(x)\;dz\\
    &= \left(\int_\Omega\Gamma^0_{ijab}(x-z)\Lambda_{ab}(z)\;dz\right)\varepsilon_{k\ell}(x)\\
    &= \left(\Gamma^0*\Lambda\right)_{ij}(x)\varepsilon_{k\ell}(x)
    \end{aligned}
\end{equation}
Putting everything together, we get
\begin{equation}
\frac{\delta \mathcal{L}}{\delta C_{ijk\ell}} = S^*_{ij}\varepsilon_{k\ell}
\qquad\text{where $S^* = S+\Gamma^0*\Lambda$.}
\label{eqn:gradC}
\end{equation}

The derivative with respect to $\overline{\varepsilon}$ only has an adjoint contribution:
\begin{equation}
    \begin{aligned}
    \frac{\delta}{\delta \overline{\varepsilon}} \int_\Omega \Lambda(z) : \mathcal{A}(\varepsilon;C,\overline{\varepsilon})\;dz &= -\int_\Omega \Lambda(z) \;dz
    \end{aligned}
\end{equation}
%%%%%%%%%%%%%%%%%%%%%%%%%%%%%%%%%%%%%%%%%%%%%%%%%%%%%%%%%%%%%%%%%%%%%%%
\subsubsection{Anisotropic case: derivatives with respect to $C(x)$}
%%%%%%%%%%%%%%%%%%%%%%%%%%%%%%%%%%%%%%%%%%%%%%%%%%%%%%%%%%%%%%%%%%%%%%%
The tensor $C$ is symmetric, i.e. $C_{ijk\ell} = C_{k\ell ij}$. This can be taken into account by setting
\begin{equation}
\frac{\delta \mathcal{L}}{\delta C_{ijk\ell}} =\begin{cases} S^*_{ij}\varepsilon_{ij} & \text{if $i=k$ and $j=\ell$}, \\
    S^*_{ij}\varepsilon_{k\ell} + \varepsilon_{ij}S^*_{k\ell}& \text{otherwise},
\end{cases}
\end{equation}
where $S^*$ is defined in \eqref{eqn:gradC}.
%%%%%%%%%%%%%%%%%%%%%%%%%%%%%%%%%%%%%%%%%%%%%%%%%%%%%%%%%%%%%%%%%%%%%%%
\subsubsection{Isotropic case: derivatives with respect to $\lambda(x)$ and $\mu(x)$}
%%%%%%%%%%%%%%%%%%%%%%%%%%%%%%%%%%%%%%%%%%%%%%%%%%%%%%%%%%%%%%%%%%%%%%%
In the isotropic case where $C_{ijk\ell} = \lambda \delta_{ij}\delta_{k\ell} + \mu (\delta_{ik}\delta_{j\ell} + \delta_{i\ell}\delta_{jk})$ we get with the chain rule
\begin{xalignat}{2}
\frac{\delta \mathcal{L}}{\delta \lambda} &= \frac{\delta \mathcal{L}}{\delta C_{ijkl}}\frac{\delta C_{ijkl}}{\delta \lambda} = \frac{\delta \mathcal{L}}{\delta C_{ijkl}} \delta_{ij}\delta_{k\ell},&
\frac{\delta \mathcal{L}}{\delta \mu} &= \frac{\delta \mathcal{L}}{\delta C_{ijkl}}\frac{\delta C_{ijkl}}{\delta \mu} = \frac{\delta \mathcal{L}}{\delta C_{ijkl}} (\delta_{ik}\delta_{j\ell}+\delta_{i\ell}\delta_{jk})
\end{xalignat}
and therefore
\begin{equation}
\begin{aligned}
\frac{\delta \mathcal{L}}{\delta \lambda} &= \operatorname{tr}\left(S^*\right)\operatorname{tr}(\varepsilon),\\
\frac{\delta \mathcal{L}}{\delta \mu} &= 2\; S^*:\varepsilon.\\
\end{aligned}
\end{equation}
with $S^*$  given in \eqref{eqn:gradC}.
%%%%%%%%%%%%%%%%%%%%%%%%%%%%%%%%%%%%%%%%%%%%%%%%%%%%%%%%%%%%%%%%%%%%%%%
\subsection{Scalar product in Voigt notation}\label{sec:voigt_scalar_product}
%%%%%%%%%%%%%%%%%%%%%%%%%%%%%%%%%%%%%%%%%%%%%%%%%%%%%%%%%%%%%%%%%%%%%%%
Care has to be taken when using Voigt notation. Representing a symmetric $3\times 3$ matrix $A$ by the six numbers
\begin{xalignat}{6}
    a_0 & = A_{00},&
    a_1 & = A_{11},&
    a_2 & = A_{22},&
    a_3 & = A_{01},&
    a_4 & = A_{02},&
    a_5 & = A_{12}
\end{xalignat}
we have for the dot-product of two matrices $A$ and $B$
\begin{equation}
A:B = a_0 b_0 + a_1 b_1 + a_2 b_2 + 2 (a_3 b_3 + a_4 b_4 + a_5 b_5) = \sum_{j=0}^{5} w_j a_j b_j =: \langle a,b\rangle_{V}.
\end{equation}
with
\begin{equation}
w^{\text{(Voigt)}} = \begin{pmatrix}1 & 1 & 1 & 2 & 2 &2\end{pmatrix}.\label{eqn:voigt_weight}
\end{equation}
This differs from the normal, Euclidean dot-product
\begin{equation}
a\cdot b = \langle a,b\rangle = \sum_{j=0}^{5} a_j b_j.
\end{equation}
In JAX the numerical values of tensors that represent derivatives (which mathematically are dual vectors) are the Riesz-representers with respect to the $\langle \cdot,\cdot \rangle$. For example, the derivative with respect to $A$ would be defined by
\begin{equation}
\frac{\delta J}{\delta A}(H) = \langle g^E, h \rangle \qquad\text{for all $h$}
\end{equation}
This differs from \eqref{eqn:riesz_representer}, where
\begin{equation}
\frac{\delta J}{\delta A}(H) = G : H = \langle g , h\rangle_V \qquad\text{for all $h$}.
\end{equation}
This implies that $g_j = g^E_j / w^{\text{(Voigt)}}_j$ with the weights $w^{\text{(Voigt)}}$ given in \eqref{eqn:voigt_weight}. To take care of this in the code we need to:
\begin{itemize}
\item Divide all incoming tensor-valued gradients elementwise by $w^{\text{(Voigt)}}$
\item Multiply all outgoing tensor-valued gradients elementwise with $w^{\text{(Voigt)}}$.
\end{itemize}
%%%%%%%%%%%%%%%%%%%%%%%%%%%%%%%%%%%%%%%%%%%%%%%%%%%%%%%%%%%%%%%%%%%%%%%
\subsection{Generic model}
%%%%%%%%%%%%%%%%%%%%%%%%%%%%%%%%%%%%%%%%%%%%%%%%%%%%%%%%%%%%%%%%%%%%%%%
While so far we assumed that $\sigma=C:\varepsilon$, we now consider a more generic model of the form $\sigma=\Sigma(\varepsilon|\theta)$ where the map $\Sigma$ depends linearly on $\varepsilon$ and $\theta$ stands for all input parameters. We assume that the differentiable function $\Sigma$ has been implemented by the user. Instead of $\delta J/\delta C(x)$ as in \eqref{eqn:func_deriv_anisotropic} or $\delta J/\delta \lambda(x)$, $\delta J/\delta \mu(x)$ in \eqref{eqn:func_deriv_isotropic} we are interested in computing the derivative
\begin{equation}
\frac{\delta J}{\delta \theta(x)}.
\end{equation}
subject to $\varepsilon(x)$ satisfying the Lippmann-Schwinger equation in \eqref{eqn:lippmann_schwinger} with
\begin{equation}
\tau = \Sigma(\varepsilon|\theta) - \lambda^0 \operatorname{tr}(\varepsilon)\mathbb{I}-2\mu^0\varepsilon\label{eqn:tau_generic}
\end{equation}
for given $\overline{\varepsilon}$. Here we have picked an isotropic, homogeneous reference material characterised by the two Lame parameters $\lambda^0$, $\mu^0$. The residual operator in \eqref{eqn:residual_operator} is replaced by
\begin{equation}
    \mathcal{A}(\varepsilon;\theta,\overline{\varepsilon}) := \varepsilon + \Gamma^0 * \left(\Sigma(\varepsilon|\theta) - \lambda^0 \operatorname{tr}(\varepsilon)-2\mu^0\varepsilon\right) - \overline{\varepsilon}\label{eqn:residual_operator}
\end{equation}
and the Lagrangian in \eqref{eqn:lagrangian} is modified appropriately. The resulting adjoint equation is
\begin{equation}
\Lambda + (\Gamma^0*\Lambda)\frac{\delta \Sigma}{\delta \varepsilon} - \lambda^0 \operatorname{tr}(\Gamma^0*\Lambda)\mathbb{I} - 2\mu^0 (\Gamma^0*\Lambda) = \mathcal{E} 
\qquad\text{with $\mathcal{E}(x) = -\left(E(x) + S(x)\frac{\delta \Sigma}{\delta \varepsilon}(x)\right)$}
\end{equation}
This results in the Lippmann-Schwinger iteration in \cref{alg:lippmann_schwinger_adjoint_generic}.
\begin{algorithm}
    \caption{Adjoint Lippmann Schwinger iteration (generic stress-strain relationship)}
    \label{alg:lippmann_schwinger_adjoint_generic}
\begin{algorithmic}[1]
    \State{Set $\Lambda=0$}
    \For{$n=0,1,\dots$}
    \State{Compute $\widehat{\Lambda}_{ij} = \mathscr{F}[\Lambda_{ij}]$}\Comment{Fourier-transform}
    \State{Set $\widehat{\Theta}_{k\ell} = -\widehat{\Gamma}^0_{k\ell ij} \widehat{\Lambda}_{ij}$}\Comment{Solve in Fourier space}
    \State{Compute $\Theta_{k\ell} = \mathscr{F}^{-1}[\widehat{\Theta}_{k\ell}]$}\Comment{Inverse Fourier-transform}        
    \State{Update $\Lambda_{ij}=\mathcal{E}_{ij} + \Theta_{k\ell}\frac{\delta \Sigma_{k\ell}}{\delta \varepsilon_{ij}} -\lambda^0 \Theta_{kk} \delta_{ij} - 2 \mu^0 \Theta_{ij}$}
    \State{Check convergence}
    \EndFor
\State{\Return $\Lambda$}
\end{algorithmic}
\end{algorithm}
$\varepsilon$ we can still be computed with \cref{alg:lippmann_schwinger_incremental}, provided $\sigma$ is obtained from $\varepsilon$ according to $\sigma=\Sigma(\varepsilon|\theta)$. Observe that the terms $S\frac{\delta \Sigma}{\delta \varepsilon}$ and $S\frac{\delta\Sigma}{\delta\varepsilon}$ are vector-Jacobian products, so they are naturally implemented as \verb~jax.vjp~'s. The derivative with respect to the parameters $\theta$ is
\begin{equation}
\left(S+\Gamma^0*\Lambda\right):\frac{\delta \Sigma}{\delta \theta}
\end{equation}
which is again a \verb~jax.jvp~.
%%%%%%%%%%%%%%%%%%%%%%%%%%%%%%%%%%%%%%%%%%%%%%%%%%%%%%%%%%%%%%%%%%%%%%%
\bibliographystyle{unsrt}
\bibliography{linear_elasticity}
\end{document}